\bookchapter{第六十七章：琢磨與發布——測試、調試與上架 App Store}

\chapterintro{App 既成，繼之何為？此章答「自能行」至「能交付」之三事：何以高效調試而覓 bug、何以書測試以防回歸、何以獻 App 於 App Store 而使天下載之。上架非徒儀式——乃使作品化為產品之分水嶺也。}

\booksection{一、調試三斧：print、斷點、預覽}

\begin{center}
\begin{tabularx}{\linewidth}{|X|X|X|}
\hline
\textbf{工具} & \textbf{司何事} & \textbf{何以用之} \tabularnewline
\hline
print + 控制臺 & 觀運行之日誌 & 碼中 print("值=\textbackslash{}(x)")，Xcode 底部控制臺觀其出 \tabularnewline
斷點調試 & 使程序止於某行，觀變量之值 & 行號旁一點而出藍標，行至彼而暫止，逐步歷之 \tabularnewline
SwiftUI 預覽 + 調試器 & 實時觀界面與視圖之層級 & \#Preview 畫布 + 運行時之 View Hierarchy 檢查器 \tabularnewline
\hline
\end{tabularx}
\end{center}

\begin{center}\small\color{gray}調試三斧\end{center}

\begin{tcbinfo}{小貼士：崩而何以讀之？}
Swift 崩時，控制臺印 “Fatal error” 及其因，Xcode 止於出誤之行而高亮。真機之舊崩，可於 Xcode Window Organizer Crashes 觀之（須付費之賬）。關鍵之詞，先搜「誤因首句 + swiftui」——問精則答精。

\end{tcbinfo}

\booksection{二、測試：為碼投保}

手測不足恃——改此一能，彼一能潛壞。測試者，悔之藥也，兼防彈之衣。iOS 用 XCTest 之框，分二種：

\begin{center}
\begin{tabularx}{\linewidth}{|X|X|X|X|}
\hline
\textbf{層級} & \textbf{測何物} & \textbf{Target} & \textbf{速度} \tabularnewline
\hline
單元測試 & 純邏輯（計算、校驗） & MyFirstAppTests & 秒級 \tabularnewline
UI 測試 & 真界面之行 & MyFirstAppUITests & 分鐘級 \tabularnewline
手動測試 & 真人之感、體驗 & —— & 恆所需也 \tabularnewline
\hline
\end{tabularx}
\end{center}

\begin{center}\small\color{gray}iOS 測試層級（模板之項已預置二 Test Target）\end{center}

\begin{lstlisting}[language=Swift]
import XCTest

// 將"金額字符串轉數字"此種邏輯抽成純函數
func parseAmount(_ input: String) -> Double? {
    guard let v = Double(input.trimmingCharacters(in: .whitespaces)),
          v > 0 else { return nil }
    return v
}

final class AmountTests: XCTestCase {
    func testEmptyInputReturnsNil() {
        XCTAssertNil(parseAmount(""))
    }
    func testNegativeReturnsNil() {
        XCTAssertNil(parseAmount("-5"))
    }
    func testNormalValueParses() {
        XCTAssertEqual(parseAmount("12.5"), 12.5)
    }
}
\end{lstlisting}

行測試：CmdU。既全綠而後改碼，心中乃有據。記賬本中金額之驗、彙總之算，皆當有單元測試護之。

\booksection{三、署名與打包：為 App 鈐「印」}

凡 iOS App 必署名而後可安——以證「此包出於汝，未遭竄改」。層級之係：

\begin{itemize}
  \item 學習之期：免費 Apple ID + Xcode 自動署名（Personal Team），真機可行，七日一續；
  \item 分發之期：付費開發者之賬（99 美元/年）簽發證書與描述之檔，而後可上架與久安。
\end{itemize}

打包：Xcode 菜單 Product Archive（歸檔），構建既成，於 Organizer 中 Distribute App 擇 App Store Connect 上傳。歸檔必以 Generic iOS Device 為的（非擬器），此新手所常忘之一程也。

\textbf{避坑之警：證書亡則更新絕}\par

署名之證與私鑰一亡，則舊 App 不得復更新——用戶惟可安「全新」之應用。善備證書（鑰匙串導出 .p12）而謹記其密碼。幸蘋果有「自動署名」與雲證書之管（Xcode 中勾 Automatically manage signing），能省手工諸坑之大半。

\booksection{四、上架 App Store 之全程}

註冊開發者賬號

準備商店素材

上傳構建包

填隱私與審覈信息

蘋果審覈（1-7 天）

上架

\begin{center}\small\color{gray}圖 5：App Store 上架六步\end{center}

\begin{itemize}
  \item 註冊：developer.apple.com，個人賬號 99 美元/年；
  \item 素材：圖標（1024×1024，勿帶圓角——蘋果自裁之）、截圖（iPhone 各尺寸）、描述、隱私政策之連結；
  \item 隱私之問卷：如實申報所收何數據——不連網之記賬本可「不收任何數據」；
  \item 審核：首提常候一至七日，見拒則依「審核之註」逐條改之——蘋果拒之由，常書之甚詳，循而行之即可；
  \item TestFlight：上架之先，以測試版予友內測（至多 10000 人），集一輪之復而後正提，專業之範十足。
\end{itemize}

\booksection{五、商店素材何以為之：第一印象即下載之量}

商店之頁，用戶「逛店」所見之櫥窗也。三素材定第一印象：

\begin{itemize}
  \item 圖標：1024×1024，勿以截圖苟且。無設計之能，則以 AI 之具生一簡潔之形，於桌面縮之而猶能一望而識，方為合格；
  \item 截圖：至少二套（iPhone 主流之尺寸），陳其最核之功能之景；每張配一句白話之文（「一秒記一筆」較「高效記賬之方」善十倍）；
  \item 描述：首句當明「此何物 + 為誰而作」，勿堆「最專、最強」——誠較誇為效。
\end{itemize}

\begin{tcbpractice}{來動手：以記賬本「交付」一過}
\booksection{六、崩潰之上報與性能優化小單}

App 上架之先，此二事宜各費三十分鐘：

\booksubsection{崩潰監控}

汝測時不崩，非謂用戶之機不崩。蘋果於 Xcode Organizer 中予免費之崩報（付費之賬可觀真機崩潰之日誌）；更省心者，接第三方 SDK（如 Firebase Crashlytics）。無崩監而上架，猶閉目而御車也。

\booksubsection{性能自察四條}

\begin{enumerate}
  \item 列表用 LazyVStack / List：長列表勿以常 VStack 硬排——懶載者，性能第一課也；
  \item 圖勿硬塞原圖：用 AsyncImage（遠圖）與系統之緩存，5MB 原圖足以滯壞舊機；
  \item 勿於主線程任重務：網絡、大檔皆行 async/await；
  \item 以 Instruments 量之：Xcode 自帶性能剖析之器（Time Profiler、記憶體洩漏之檢）——先量而後優，勿憑感覺。
\end{enumerate}

(1) 為 parseAmount 等邏輯補三單元測試（CmdU 行之而綠）；(2) 以斷點親捕自造之崩而修之；(3) Product Archive 打一包，由 TestFlight 邀友內測；(4) 有餘力則註冊開發者之賬，以記賬本真提而上架。前三程今夕可成。

\end{tcbpractice}

\begin{tcbsummary}{本章小結}
\begin{itemize}
  \item 調試三斧：print 之誌、斷點、SwiftUI 預覽與視圖檢查器
  \item XCTest：單元測試（CmdU）+ UI 測試 + 手動測試，模板已預置
  \item 署名：免費賬號自動署名（七日一續）；付費賬號方能分發
  \item 打包：Generic iOS Device + Archive + Distribute；證書務必備份
  \item 上架六步：註冊 素材 上傳 隱私問卷 審核 上架；TestFlight 先內測
\end{itemize}

\end{tcbsummary}
